\documentclass[10pt,twocolumn]{article}
\usepackage[margin=0.8in]{geometry}
\usepackage{booktabs}
\usepackage[hidelinks]{hyperref}
\usepackage{url}

\title{Bind: A Registry of Legal--Architectural Incompatibilities}
\author{Will Foti}
\date{November 2026}

\begin{document}
\maketitle

\begin{abstract}
We publish 40 cases in which a legal requirement and a technical architecture
issue contradictory mandates, drawn from 7 jurisdictions and 110 cited sources.
A case enters the registry only if a competent engineer with unlimited budget
still cannot satisfy both sides; 11 candidates failed that test and are
published as rejections. We group the cases into six shapes. An adversarial
prior-art audit by a non-Claude model found that four of the six shapes are
established ideas under other names, so we do not claim the shapes as new. The
contribution is the case set, the inclusion test, and one empirical result:
of the 35 cases rated unresolvable, five have stopped being live, and none of
the five was ended by engineering. Three ended through legislation, one through
a court, and one through abandonment of the product.
\end{abstract}

\section{Inclusion test}

A case qualifies only if a competent engineer with unlimited budget still
cannot satisfy both mandates. Cases that are merely expensive, or that a better
implementation would fix, are out. Each case carries one of three
resolvability ratings: \emph{unresolvable} (35 cases), \emph{resolvable at
cost} (3), or \emph{technology gap}, an open research problem (2). Cases
rated resolvable at cost or technology gap stay in the registry as boundary
cases, because they mark where the test stops.

Inclusion verdicts are recorded per case: 29 pass, 4 fail and are kept as
boundary cases, and 7 are undecided. The 11 rejected candidates are published
with a reason: resolved by a known technique (6), merely expensive (3), or out
of scope (2). HIPAA versus cloud hosting is a typical rejection. Business
associate agreements settled it a decade ago.

\section{Shapes}

Table~\ref{tab:shapes} lists the six active shapes. A seventh, Locality
vs.\ Emergent Placement, was retired in taxonomy v2 because every one of its
cases could be resolved by pinning data in-country. Shape assignments are
versioned, so its one case (C1) keeps that history.

Each shape names an owner, the party who could end the conflict:
the architect, the legislator, a researcher, or nobody. Sovereign
Contradiction has no owner. Two sovereigns issue opposite commands about one
artifact, and the engineer holding the data cannot satisfy both.

\begin{table*}[t]
\centering
\caption{Shapes of legal--architectural incompatibility, with case counts and prior-art verdicts.}
\label{tab:shapes}
\begin{tabular}{rllrl}
\toprule
\# & Shape & Family & Cases & Prior-art verdict \\
\midrule
  1 & Mandate--Mutation & architectural & 9 & partial \\
  2 & Sovereign Contradiction & conflict of law & 9 & rebrand \\
  4 & Deliberate Blindness & architectural & 9 & rebrand \\
  5 & Trust-Root Monopoly vs. Plurality & architectural & 6 & partial \\
  6 & Verification vs. Minimization & research gap & 8 & rebrand \\
  7 & Binary vs. Continuum & drafting error & 10 & rebrand \\
\bottomrule
\end{tabular}
\end{table*}

\subsection{Prior-art audit}

Each shape was submitted to gpt-5 with a prompt instructing it to find prior
work describing the same mechanism and to return \emph{novel},
\emph{partial}, \emph{rebrand}, or \emph{inconclusive}. The prompt is
\path{docs/prompts/ADVERSARIAL_PRIOR_ART_PROMPT.md}, version 1. A Claude model
drafted the shapes, so Claude models were excluded as auditors.

\begin{itemize}
\item \textbf{Sovereign Contradiction: rebrand.} This is the ``true
  conflict'' of conflict-of-laws doctrine~\cite{hartfordfire,pauwelyn2003}.
\item \textbf{Deliberate Blindness: rebrand.} The conflict between lawful
  access and end-to-end encryption is described in RFC~1984, RFC~2804, and
  \emph{Keys Under Doormats}~\cite{rfc1984,rfc2804,abelson2015keys}. Our cases
  extend it to no-log VPNs, confidential computing, and breach notification.
\item \textbf{Verification vs.\ Minimization: rebrand.} Chaum and Brands name
  it and propose selective disclosure as the
  fix~\cite{chaum1985,brands2000}.
\item \textbf{Binary vs.\ Continuum: rebrand.} Ohm and Schwartz and Solove
  argue that the PII threshold sits on a continuum~\cite{ohm2010,schwartzsolove2011};
  Endicott treats bright lines over continua in law
  generally~\cite{endicott2000}.
\item \textbf{Mandate--Mutation: partial.} Covered domain by domain (erasure
  vs.\ ledgers, erasure vs.\ model weights~\cite{caoyang2015}), without a
  single name across domains.
\item \textbf{Trust-Root Monopoly vs.\ Plurality: partial.} The single-root
  argument appears for DNS~\cite{rfc2826} and the Web PKI~\cite{clark2013sok}.
  The audit found no account that applies it to repair and sideloading
  mandates.
\end{itemize}

We kept all four rebranded shapes. A rebrand verdict says the mechanism has
prior art. It says nothing about whether the cases are correctly classified,
and retiring the four shapes would leave 24 of the 39 shaped cases with no shape. Each
shape's provenance note now cites its prior art.

\section{Corpus}

Table~\ref{tab:corpus} lists all 40 cases. Each case records the mandating
instruments with pinpoint provisions, the architectural property in conflict,
a resolvability rating, and, where the conflict ended, the mechanism and the
instrument that ended it.

\begin{table*}[t]
\centering
\caption{Corpus of documented incompatibilities.}
\label{tab:corpus}
\begin{tabular}{llll}
\toprule
Ref & Case & Shape & Resolvability \\
\midrule
  A1 & GDPR Art. 17 erasure vs. append-only ledgers & Mandate--Mutation & resolvable costly \\
  A2 & Erasure/deletion rights vs. mandatory retention & Mandate--Mutation & unresolvable \\
  A3 & Erasure vs. transparency logs & Mandate--Mutation & unresolvable \\
  A4 & Erasure vs. trained model weights & Mandate--Mutation & unresolvable \\
  B1 & Schrems II transfer restrictions vs. CLOUD Act production ord… & Sovereign Contradiction & unresolvable \\
  B2 & eIDAS 2.0 QWAC trust-root mandate vs. root program distrust a… & Sovereign Contradiction & unresolvable \\
  C1 & Data localization vs. anycast + global load balancing &  & resolvable costly \\
  D1 & E2EE vs. Art. 33/34 breach notification & Deliberate Blindness & unresolvable \\
  D2 & UK Online Safety Act s.121 scanning notices vs. E2EE & Deliberate Blindness & unresolvable \\
  D3 & Exposure notification: DP-3T/GAEN vs. epidemiological need & Deliberate Blindness & unresolvable \\
  E1 & Right to repair vs. parts pairing / attested firmware & Trust-Root Monopoly vs. Plurality & unresolvable \\
  E2 & DMA Art. 6(4) sideloading vs. single-root iOS distribution mo… & Trust-Root Monopoly vs. Plurality & unresolvable \\
  E3 & DMA Art. 7 messaging interoperability vs. E2EE guarantees & Trust-Root Monopoly vs. Plurality & unresolvable \\
  E4 & Client attestation (Play Integrity, WEI) vs. accessibility an… & Trust-Root Monopoly vs. Plurality & unresolvable \\
  F1 & Age verification mandates vs. GDPR Art. 5(1)(c) & Verification vs. Minimization & tech gap \\
  F2 & FATF Travel Rule vs. self-hosted wallets & Verification vs. Minimization & unresolvable \\
  G1 & "Anonymous" as a binary vs. re-identifiability as a continuum & Binary vs. Continuum & unresolvable \\
  G2 & Art. 16 rectification vs. probabilistic inference & Binary vs. Continuum & unresolvable \\
  G3 & Art. 15 / Art. 22 explainability vs. model opacity + trade se… & Binary vs. Continuum & unresolvable \\
  G4 & Differential privacy budget vs. redistricting accuracy & Binary vs. Continuum & unresolvable \\
  G5 & AI Act Art. 10(5) bias testing vs. GDPR Art. 9 & Binary vs. Continuum & unresolvable \\
  H1 & EU data retention mandates vs. ePrivacy Art. 15(1) & Sovereign Contradiction & unresolvable \\
  H2 & PCI DSS 10.x / SOC2 logging vs. deletion rights over IPs and… & Mandate--Mutation & unresolvable \\
  H3 & FRCP 37(e) litigation hold vs. ephemeral messaging & Mandate--Mutation & unresolvable \\
  N4 & Sanctions Blocking vs. Immutable Contracts & Mandate--Mutation & unresolvable \\
  N5 & EU Blocking Statute vs. OFAC Iran Rules & Sovereign Contradiction & unresolvable \\
  N6 & US Correspondent Subpoenas vs. China Data Export & Sovereign Contradiction & unresolvable \\
  N7 & MiCA Scope vs. Partial Decentralisation & Binary vs. Continuum & unresolvable \\
  N8 & Sensitive Care Segregation vs. Clinical Records & Binary vs. Continuum & unresolvable \\
  N19 & EDR Anonymity vs. Crash Trace & Binary vs. Continuum & unresolvable \\
  N20 & Stablecoin Burn Orders vs. MiCA Redemption & Sovereign Contradiction & unresolvable \\
  N12 & UK IPA technical capability notices vs. provider-blind E2EE b… & Deliberate Blindness & unresolvable \\
  N13 & India IT Rules 4(2) first-originator traceability vs. E2EE me… & Verification vs. Minimization & unresolvable \\
  N14 & DMA 6(7) Wi-Fi interoperability vs. single-account trust doma… & Trust-Root Monopoly vs. Plurality & unresolvable \\
  N15 & DSA monthly active recipients vs. no individual tracking & Verification vs. Minimization & unresolvable \\
  N16 & DMA 6(8) independent ad verification vs. data minimization & Verification vs. Minimization & tech gap \\
  N17 & DSA 26(3) special-category ad ban vs. behavioural profiling & Binary vs. Continuum & unresolvable \\
  N18 & Russia 149-FZ decoding-key duty vs. E2EE messaging & Deliberate Blindness & unresolvable \\
  N21 & CERT-In VPN subscriber records vs. no-log VPN & Deliberate Blindness & unresolvable \\
  N22 & DSA Art. 28 age assurance robustness vs. no additional data & Verification vs. Minimization & resolvable costly \\
\bottomrule
\end{tabular}
\end{table*}

\section{Resolution}

Six cases have stopped being live.

\begin{center}
\begin{tabular}{lll}
\toprule
Ref & Rating & Ended by \\
\midrule
B2 & unresolvable & Legislation (eIDAS 2.0 Art.~45a(2)) \\
G5 & unresolvable & Legislation (AI Act Art.~10(5)) \\
N15 & unresolvable & Legislation (DSA recital 77) \\
N4 & unresolvable & Court (\emph{Van Loon}, 5th Cir.\ 2024) \\
N21 & unresolvable & Product abandonment \\
N22 & resolvable at cost & Engineering \\
\bottomrule
\end{tabular}
\end{center}

The only engineering resolution is N22, age assurance under DSA Art.~28,
which the registry had already rated resolvable at cost. Its fix is
double-blind age tokens with device-local holder binding. No case rated
unresolvable has been ended by engineering. This is consistent with the
inclusion test, and it is the claim the registry will test as cases
accumulate. An unresolvable case resolved by engineering would show a rating
was wrong.

\section{Method and limitations}

\begin{itemize}
\item \textbf{Drafting.} The first 24 cases come from a working corpus
  (v0.1) whose citations were recalled from memory. Claude Opus 5.5 drafted
  the remaining 16 cases and checked every citation against a primary source.
  Each check is recorded in a verification ledger with its URL.
\item \textbf{Review.} The author approved the corpus for publication on the
  strength of the model's source verification, without an independent
  case-by-case human check. Every case records this sign-off.
\item \textbf{Unverified source.} One citation is unverified:
  \emph{WhatsApp v.\ Union of India} (Delhi HC, filed 2021). Secondary
  reporting confirms the filing. We found no judgment, and the court's case
  status page could not be queried automatically.
\item \textbf{Single auditor.} One model ran the prior-art audit once. A
  second non-Claude auditor would show whether the verdicts are stable.
\item \textbf{Open verdicts.} Seven cases have an undecided inclusion
  verdict.
\end{itemize}

\bibliographystyle{plain}
\bibliography{priorart,bind}

\end{document}
